% VLI* theoretical section. Compile standalone with pdflatex if desired.
\documentclass[11pt]{article}
\usepackage[margin=1in]{geometry}
\usepackage{amsmath,amssymb,amsthm,mathtools}
\usepackage[T1]{fontenc}
\usepackage{microtype}
\newtheorem{assumption}{Assumption}
\newtheorem{lemma}{Lemma}
\newtheorem{theorem}{Theorem}
\newtheorem{proposition}{Proposition}
\newtheorem{corollary}{Corollary}
\newcommand{\R}{\mathbb{R}}
\newcommand{\one}{\mathbf{1}}
\newcommand{\pos}[1]{[#1]_+}

\begin{document}
\section{Bias-gradient label recovery and VLI$^*$}
\label{sec:vli-star}

\paragraph{Scope.} This section studies the class counts in one client window
from the accumulated cross-entropy gradient of the final-layer bias. It does
not claim that the difficulty vector or confusion matrix introduced below can
be recovered from a shared update. Their estimation is a separate problem.
The main results are: (i) the exact relationship between bias gradients and
per-class error mass; (ii) a precise sense in which auxiliary-only
calibrations lose label dependence as training error vanishes; and (iii) an
identifiable estimator when class difficulty and the error-confusion geometry
are known.

\subsection{Setup and notation}
Let a client window contain $n$ examples from $C$ classes. For example $i$,
let $y_i\in\{1,\ldots,C\}$ be its label, $z_i\in\R^C$ the logits, and
$p_i=\operatorname{softmax}(z_i)$. The cross-entropy loss is
$\ell_i=-\log p_{i,y_i}$. Define the bias-gradient signal
\[
  G=\sum_{i=1}^n(p_i-e_{y_i}), \qquad
  N_j=\#\{i:y_i=j\}, \qquad \sum_jN_j=n.
\]
The code accumulates mean-loss gradients per minibatch. When all minibatches
have size $B$, that signal differs from the example-sum convention above by a
common positive factor (and possibly by a further common window normalization).
All scale-invariant formulas below are unchanged by such a factor. A server
weight update has the opposite sign under plain SGD; here $G$ always denotes
the loss gradient, not the signed parameter update.

For each example let $\epsilon_i=1-p_{i,y_i}$ and
$\bar\epsilon=n^{-1}\sum_i\epsilon_i$. For a present class $j$, let
\[
 \bar\epsilon_j=\frac{1}{N_j}\sum_{i:y_i=j}\epsilon_i,
 \qquad a_j=N_j\bar\epsilon_j.
\]
For absent classes set $a_j=0$; $\bar\epsilon_j$ is undefined there and a
positive placeholder may be used for $r_j$ because it is multiplied by
$a_j=0$. Define the off-diagonal error-transition matrix
\[
 Q_{jk}=\begin{cases}
 \displaystyle\frac{\sum_{i:y_i=j}p_{i,k}}{a_j},& k\ne j,\ a_j>0,\\
 0,& k=j.
 \end{cases}
\]
Every defined row is nonnegative and sums to one. When a class has zero error
mass, its row can be assigned any probability distribution over the other
classes; it contributes nothing to $Q^Ta$. Define
\[
 \bar\epsilon=\frac1n\sum_j a_j,\qquad
 r_j=\frac{\bar\epsilon_j}{\bar\epsilon}\quad(N_j>0).
\]
Thus $a_j=N_jr_j\bar\epsilon$ for present classes.

\subsection{Exact gradient identities}
\begin{lemma}[Gradient norm and training loss]
\label{lem:gradient-bound}
For every example, $\|p_i-e_{y_i}\|_1=2\epsilon_i$. Consequently,
\[
 \|G\|_1\le 2n\bar\epsilon\le 2nL_{\rm train},
 \qquad L_{\rm train}=\frac1n\sum_i\ell_i.
\]
\end{lemma}
\begin{proof}
The true-class coordinate has absolute value $1-p_{i,y_i}=\epsilon_i$.
All other coordinates are nonnegative and sum to $1-p_{i,y_i}=\epsilon_i$.
This proves the per-example equality. The triangle inequality gives
$\|G\|_1\le\sum_i\|p_i-e_{y_i}\|_1=2n\bar\epsilon$.
For $x\in(0,1]$, $1-x\le-\log x$; apply it to $x=p_{i,y_i}$ and average.
\end{proof}

\begin{lemma}[Count, error-mass, and confusion identities]
\label{lem:exact-identities}
Let $P_j=\sum_i p_{i,j}$. Then
\[
 G_j=P_j-N_j=-a_j+E_j,\qquad
 E_j=\sum_{i:y_i\ne j}p_{i,j},\qquad
 N_j=P_j-G_j.
\]
Moreover, for any row-stochastic $Q$ defined above,
\[
                     G=-(I-Q^T)a. \tag{1}
\]
If $N_j=0$, then $G_j=P_j\ge0$.
\end{lemma}
\begin{proof}
The $j$th coordinate of $G$ sums all predicted probability assigned to class
$j$ and subtracts one for every true class-$j$ example, yielding $P_j-N_j$.
Separate $P_j$ into its true-class contribution
$\sum_{i:y_i=j}p_{i,j}=N_j-a_j$ and its other-class contribution $E_j$.
This gives the first identities and the absent-class sign property. For each
target class $k$, the incoming error mass is
\[
 E_k=\sum_{j\ne k}\sum_{i:y_i=j}p_{i,k}
     =\sum_j a_jQ_{j,k}=(Q^Ta)_k.
\]
Substituting in $G=-a+E$ proves (1).
\end{proof}

\begin{theorem}[Auxiliary-only calibration becomes label-independent]
\label{thm:aux-collapse}
Let an estimator of normalized class proportions have the form
\[
 \widehat\pi=\Phi(c-G/n),
\]
where $c$ depends only on information outside the attacked client window and
$\Phi$ is $L_\Phi$-Lipschitz in a stated norm. Then
\[
 \|\widehat\pi-\Phi(c)\|
 \le \frac{2L_\Phi}{n}\|G\|_1
 \le 4L_\Phi\bar\epsilon
 \le 4L_\Phi L_{\rm train}. \tag{2}
\]
Hence as $L_{\rm train}\to0$, the estimate approaches a value determined by
$c$, not by the attacked labels. If an estimator returns counts $n\Phi(\cdot)$,
the corresponding count bound is multiplied by $n$.
\end{theorem}
\begin{proof}
The Lipschitz property gives
$\|\Phi(c-G/n)-\Phi(c)\|\le L_\Phi\|G/n\|_1$. Apply
Lemma~\ref{lem:gradient-bound}. The limit follows directly.
\end{proof}
\paragraph{Interpretation and limitation.} The theorem applies to an
implementation only after its exact map, normalization, clipping, and
rounding have been shown to satisfy this representation and Lipschitz
condition. In particular, it is not by itself a proof that every RLU posterior
search or every auxiliary-calibrated heuristic in the repository satisfies
(2). Simplex normalization can also fail to be globally Lipschitz near a
zero denominator.

\begin{proposition}[Held-out error gap]
\label{prop:heldout-gap}
For any example misclassified by the argmax rule, $p_y\le1/2$ and therefore
$1-p_y\ge1/2$. If the held-out error rate is $\tau_{\rm err}$, then
\[
 \bar\epsilon_{\rm aux}\ge\tau_{\rm err}/2.
\]
If simultaneously $\bar\epsilon_{\rm train}\to0$ and
$\tau_{\rm err}$ stays bounded away from zero, then
$\bar\epsilon_{\rm aux}/\bar\epsilon_{\rm train}\to\infty$.
\end{proposition}
\begin{proof}
On a misclassified example some $k\ne y$ has $p_k\ge p_y$. Since
$p_k+p_y\le1$, $p_y\le1/2$. Average the resulting lower bound on
$1-p_y$ over the held-out set. The ratio statement follows.
\end{proof}

\subsection{Identifiability from known error geometry}
\begin{assumption}[Known irreducible error geometry]
\label{ass:geometry}
The matrix $Q$ is known, row-stochastic, nonnegative, has zero diagonal, and
is irreducible. The relative difficulty vector $r$ is known and strictly
positive on present classes. At least one class is absent, and $\bar\epsilon>0$.
\end{assumption}

\begin{theorem}[Recovery with known $(r,Q)$]
\label{thm:general-recovery}
Under Assumption~\ref{ass:geometry}, the nonnegative solution of (1) whose
smallest feasible point lies on the boundary determines $a$ uniquely. If
$v>0$ is the normalized Perron vector satisfying $Q^Tv=v$ and
$\one^Tv=1$, then every solution of (1) is
\[
 a=a_0+tv,\qquad a_0=(I-Q^T)^+(-G).
\]
The condition $a\ge0$ and at least one absent class selects
\[
 t=\max_j\left(-\frac{a_{0,j}}{v_j}\right),
 \qquad
 \widehat N_j=n\frac{a_j/r_j}{\sum_k a_k/r_k}. \tag{3}
\]
For exact $G,r,Q$, (3) equals the true class counts.
\end{theorem}
\begin{proof}
Since $Q$ is irreducible and stochastic, Perron--Frobenius theory gives a
one-dimensional nullspace of $I-Q^T$, spanned by a strictly positive $v$.
Thus all solutions of (1) differ by $tv$. The constraint
$a_{0,j}+tv_j\ge0$ for every $j$ is equivalent to
$t\ge\max_j(-a_{0,j}/v_j)$. At this smallest feasible $t$, at least one
coordinate is zero. Conversely, the true $a$ is nonnegative and has a zero
coordinate because at least one class is absent; it must be this boundary
solution. For each present class, $a_j=N_jr_j\bar\epsilon$. Therefore
$a_j/r_j=N_j\bar\epsilon$, and normalizing these values to sum to $n$ gives
the true counts. Absent-class terms are zero regardless of the positive
placeholder used for their $r_j$.
\end{proof}
\paragraph{Noisy observations.} With noisy $G$ or misspecified $Q$, the
linear system may not have an exact nonnegative solution. A practical method
must report the residual and use a constrained least-squares or regularized
fit. Reducible $Q$ can have a higher-dimensional nullspace and is not covered
by this identifiability theorem.

\begin{corollary}[Difficulty-mismatch bound]
\label{cor:kappa}
Suppose $Q$ is known and the boundary solution $a$ is exact, but the estimator
uses uniform difficulty $r_j=1$ rather than the true $r$. Let
$\kappa=r_{\max}/r_{\min}$ over present classes. Then its instance recovery
rate, defined as $\sum_j\min(\widehat\pi_j,\pi_j)$, is at least
$1/\kappa$.
\end{corollary}
\begin{proof}
The estimate is $q_j=a_j/A$, where $A=\sum_k a_k$. The true proportion is
$\pi_j=(a_j/r_j)/Z$, where $Z=\sum_k a_k/r_k$. Hence
\[
 \frac{q_j}{\pi_j}=r_j\frac{Z}{A}.
\]
The factor $Z/A$ is a weighted average of $1/r_k$, so it is at least
$1/r_{\max}$. Thus $q_j\ge\pi_j/\kappa$ for every present class, and
$\min(q_j,\pi_j)\ge\pi_j/\kappa$. Summing proves the claim.
\end{proof}

\subsection{Symmetric closed form and relation to VLI peeling}
\begin{assumption}[Symmetric error geometry]
\label{ass:symmetric}
All present classes have the same mean error, so $r_j=1$. Every class
distributes its error mass uniformly over the other $C-1$ classes:
$Q_{jk}=1/(C-1)$ for $k\ne j$.
\end{assumption}

\begin{theorem}[Closed-form VLI$^*$]
\label{thm:closed-form}
Under Assumption~\ref{ass:symmetric} and at least one absent class,
\[
 G_j=\frac{n\bar\epsilon-Ca_j}{C-1},\qquad
 G_{\max}=\frac{n\bar\epsilon}{C-1},\qquad
 N_j=\frac{n}{C}\left(1-\frac{G_j}{G_{\max}}\right). \tag{4}
\]
The formula is invariant to multiplying $G$ by a positive scalar.
\end{theorem}
\begin{proof}
Under uniform confusion, $(Q^Ta)_j=(\sum_k a_k-a_j)/(C-1)$. By (1),
\[
 G_j=-a_j+\frac{\sum_k a_k-a_j}{C-1}
     =\frac{\sum_k a_k-Ca_j}{C-1}.
\]
Since $\sum_k a_k=n\bar\epsilon$, this proves the first equality. An absent
class has $a_j=0$, so it attains $G_{\max}=n\bar\epsilon/(C-1)$; no other
coordinate can exceed it because $a_j\ge0$. Rearranging and using
$a_j=N_j\bar\epsilon$ gives (4). A common positive scale cancels in the
ratio $G_j/G_{\max}$.
\end{proof}

Let $S_- = \sum_{j:G_j<0}G_j<0$ and let the synthetic basis have diagonal
$d<0$ and off-diagonal $-d/(C-1)$. The signal-matched diagonal magnitude is
\[
 d^*=-\frac{(C-1)G_{\max}}{n},\qquad
 \operatorname{boost}^*=\frac{d^*}{S_-/n}
 =\frac{(C-1)G_{\max}}{-S_-}. \tag{5}
\]
Here the code's base diagonal is $S_-/n$, so its boost is positive. (Equation
(5) writes the ratio using magnitudes; equivalently
$\operatorname{boost}^*=(C-1)G_{\max}/\sum_j(-G_j)_+$.) With this diagonal,
each class-count increment subtracts exactly its symmetric per-example basis
vector, up to the common signal scale. Greedy peeling therefore has the same
residual scores as the closed-form proportions in (4), and integer
largest-remainder rounding gives the same counts when the model is exactly
on the symmetric model.

If only one class is present, $G_{\max}=n\bar\epsilon/(C-1)$ and
$-S_-=n\bar\epsilon$, so $\operatorname{boost}^*=1$. In the symmetric case
with several present classes, the ratio adjusts the basis magnitude to the
error scale inferred from an absent class. This replaces the current
scale-dependent heuristic factors with a directly derived value.

\subsection{Margin averaging caveat}
For a sample with margin
$M=z_y-\log\sum_{k\ne y}e^{z_k}$, its error is exactly
$\epsilon=(1+e^M)^{-1}$. When $M$ is large and positive,
$\epsilon\approx e^{-M}$. If margins within a class are approximately
$\mathcal N(\mu,\sigma^2)$ and the large-margin approximation is valid over
most of their mass, then
\[
 \mathbb E[e^{-M}]=e^{-\mu+\sigma^2/2},
 \qquad e^{-\mathbb E[M]}=e^{-\mu}.
\]
Thus plugging in a mean margin underestimates this exponential surrogate by
$e^{\sigma^2/2}$. Jensen's inequality gives
$\mathbb E[e^{-M}]\ge e^{-\mathbb E[M]}$. This is an asymptotic explanation,
not an exact formula for logistic error: a Gaussian margin has nonzero mass
outside the large-positive-margin regime, where $\epsilon\approx e^{-M}$ is
not accurate.

\subsection{What the theorems do not establish}
The exact identities hold for per-example cross-entropy bias gradients,
including examples evaluated at different local model states. They do not
automatically identify a server's full optimizer update when it also contains
momentum state, proximal terms, weight decay, or FedDyn correction terms.
The current generalized implementation takes $(r,Q)$ as inputs; these are
oracle diagnostics in the accompanying evaluation script. Recovering them
from public model weights, held-out auxiliary data, or the shared update is
an open step. Finally, if every class appears in the attacked window, the
absent-class boundary used in (3)--(4) is unavailable.

\end{document}
